\documentclass{article}


% if you need to pass options to natbib, use, e.g.:
%     \PassOptionsToPackage{numbers, compress}{natbib}
% before loading neurips_2024


% ready for submission
\usepackage{neurips_2024}


% to compile a preprint version, e.g., for submission to arXiv, add add the
% [preprint] option:
%     \usepackage[preprint]{neurips_2024}


% to compile a camera-ready version, add the [final] option, e.g.:
%     \usepackage[final]{neurips_2024}


% to avoid loading the natbib package, add option nonatbib:
%    \usepackage[nonatbib]{neurips_2024}


\usepackage[utf8]{inputenc} % allow utf-8 input
\usepackage[T1]{fontenc}    % use 8-bit T1 fonts
\usepackage{hyperref}       % hyperlinks
\usepackage{url}            % simple URL typesetting
\usepackage{booktabs}       % professional-quality tables
\usepackage{amsfonts}       % blackboard math symbols
\usepackage{nicefrac}       % compact symbols for 1/2, etc.
\usepackage{microtype}      % microtypography
\usepackage{xcolor}         % colors


\title{Preliminary Report For Project 7\\Sentiment Analysis}


% The \author macro works with any number of authors. There are two commands
% used to separate the names and addresses of multiple authors: \And and \AND.
%
% Using \And between authors leaves it to LaTeX to determine where to break the
% lines. Using \AND forces a line break at that point. So, if LaTeX puts 3 of 4
% authors names on the first line, and the last on the second line, try using
% \AND instead of \And before the third author name.


\author{
  Nolan Capomaggio – Mathis Morra-Fisher\\
  Group 19\\
  Warsaw University Of technologies\\
}


\begin{document}


\maketitle

\begin{abstract}
  Sentiment analysis is a key task in Natural Language Processing (NLP), 
  aiming to automatically determine the emotional polarity of a given text. 
  In this work, the goal is to address binary sentiment classification of Amazon 
  product reviews. We distinguish positive (4–5 stars) from negative (1–2 stars) 
  review. We use the Amazon Reviews dataset (Kaggle, Adame Bittlingmayer), 
  which includes approximately 4 million labeled reviews. We propose to compare 
  three approaches of increasing complexity: Logistic Regression with TF-IDF 
  features, a Bi-LSTM recurrent neural network, and a fine-tuned DistilBERT 
  transformer model. Models are evaluated using accuracy and macro F1-score 
  on a held-out test set.
\end{abstract}


\section{Introduction}



\section{Dataset}


\subsection{Description}
\subsection{Data Balance \& Splits}
\subsection{Preprocessing}
\subsection{Evaluation Metrics}

\section{Solution Description}

\section{Conclusion}


\section*{References}

{\small
[1] NeurIPS\ (2024) Call for Papers\\ \url{https://neurips.cc/Conferences/2024/CallForPapers}

[2] Adam Bittlingmayer\ (2020) Amazon Reviews for Sentiment Analysis\\ \url{https://www.kaggle.com/datasets/bittlingmayer/amazonreviews}
}


%%%%%%%%%%%%%%%%%%%%%%%%%%%%%%%%%%%%%%%%%%%%%%%%%%%%%%%%%%%%

\appendix

\section{Appendix / supplemental material}


Optionally include supplemental material (complete proofs, additional experiments and plots) in appendix.
All such materials \textbf{SHOULD be included in the main submission.}

\end{document}